% sections/position.tex -- artistic-practitioner position (shared).
% Facts from the author's CV and publications only; see source_verification.md, F.
I come to this question as a practitioner. Architecture at Tabriz Islamic
Art University and landscape architecture at Istanbul Technical University
taught me to treat designed environments as media of social life:
thresholds, paths and light shape who gathers, who is seen and how people
move together. In ITU's Virtual Landscape Lab I studied how people perceive
mediated environments \citep{jahandideh2023}. My current master's in
Interactive Arts and Technology at Simon Fraser University, in Steve
DiPaola's iViz Lab, moved this practice into interactive 3D and game
environments: I build multiplayer worlds and physiological interfaces in
Unity and C\#, use image-generation models in design work, and study these
systems with participants. One study is accepted at ACM ISS 2026
\citep{jahandideh_teammate_ms}; another is under revision
\citep{jahandideh_bioterrarium_ms}. The through-line runs from designing
spaces, to designing mediated environments, to asking who controls how a
technical system represents and translates people: in the teammate study,
participants accepted a shared heart-rate cue only if they controlled who saw
it. As artist, system designer, facilitator, participant observer and analyst,
I will hold power over what participants can do and how their actions are
read, so documenting my own decisions is part of the method.
